\documentclass[11pt]{article}
\usepackage[margin=1in]{geometry}
\usepackage{booktabs}
\usepackage{longtable}
\usepackage{array}
\title{P12B Macro Project: Forward Brainstorm (Option Map, Not Fixed Plan)}
\author{Steven Chung}
\date{\today}

\begin{document}
\maketitle

\section{Decision Rule (Hard Filter)}
Continue only if an option satisfies both:
\begin{enumerate}
\item Strong portfolio signal for applied economist / economic data analyst roles.
\item Realistic delivery in 14 days or less with stable execution.
\end{enumerate}

If either fails, stop or defer that branch.

\section{Branch Map: All Meaningful Directions}

\subsection{A. Identification and Causal Inference Branches}
\textbf{A1. Country and year fixed effects (two-way FE)}
\begin{itemize}
\item Goal: absorb time-invariant country heterogeneity and global shocks.
\item Value: standard baseline for panel credibility.
\item Risk: still vulnerable to time-varying omitted variables and simultaneity.
\end{itemize}

\textbf{A2. FE plus controls}
\begin{itemize}
\item Add controls such as trade openness, fiscal proxy, financial depth, exchange-rate regime proxy if available.
\item Value: better model defensibility in interviews.
\item Risk: data collection burden and model instability if controls are weak/incomplete.
\end{itemize}

\textbf{A3. IV / 2SLS branch}
\begin{itemize}
\item Candidate idea set: external monetary pressure proxies, lagged external instruments, commodity-linked external shocks for suitable subsamples.
\item Mandatory diagnostics: first-stage F, weak-IV robust inference, over-identification checks (if over-identified).
\item Value: highest signal if instrument story is defensible.
\item Risk: instrument validity is the main failure point.
\end{itemize}

\textbf{A4. Dynamic panel branch}
\begin{itemize}
\item Add lagged dependent variables and dynamic adjustment story.
\item Value: closer to macro transmission logic.
\item Risk: complexity rises quickly; easy to over-scope beyond 14 days.
\end{itemize}

\subsection{B. Robustness and Stress-Test Branches}
\textbf{B1. Alternative variable definitions}
\begin{itemize}
\item Money proxy variants (M2 level growth vs ratio transformations).
\item Inflation proxy variants (CPI-based alternatives where feasible).
\item Value: demonstrates that findings are not a single-definition artifact.
\end{itemize}

\textbf{B2. Sample slicing}
\begin{itemize}
\item By period (pre/post major global events), income groups, inflation regimes.
\item Value: strengthens external validity narrative.
\item Risk: smaller cells reduce statistical power.
\end{itemize}

\textbf{B3. Outlier influence checks}
\begin{itemize}
\item Winsorization / robust regression / leave-one-country-out checks.
\item Value: protects against ``single-country drives everything'' critique.
\end{itemize}

\subsection{C. Portfolio Productization Branches}
\textbf{C1. Reproducible research package}
\begin{itemize}
\item One canonical notebook, one exported dataset, one methods note, one results note.
\item Value: immediate hiring signal for execution quality.
\end{itemize}

\textbf{C2. Interview-facing one-pager}
\begin{itemize}
\item Problem, method choice, key result, caveat, business relevance.
\item Value: high communication ROI.
\end{itemize}

\textbf{C3. Light dashboard/static chart sheet}
\begin{itemize}
\item Non-interactive static figures acceptable if consistent and clean.
\item Value: faster than full app build; avoids over-engineering.
\end{itemize}

\subsection{D. Adjacent Extension Branches}
\textbf{D1. Forecasting extension}
\begin{itemize}
\item Use money-growth information as one feature in inflation forecasting benchmarks.
\item Value: links causal and predictive framing.
\item Risk: can dilute core identification story if too broad.
\end{itemize}

\textbf{D2. Policy / risk memo variant}
\begin{itemize}
\item Convert findings into a short macro-risk interpretation memo.
\item Value: direct relevance for risk and analytics roles.
\end{itemize}

\section{Portfolio Value Matrix (Initial Scoring)}
Scale: High / Medium / Low.

\begin{center}
\begin{tabular}{p{4.0cm}p{2.0cm}p{2.0cm}p{2.2cm}p{3.2cm}}
\toprule
Option & Portfolio Signal & Causal Credibility & 14-Day Feasibility & Recommended Status \\
\midrule
Two-way FE baseline & High & Medium & High & Do first if continuing \\
FE + controls & High & Medium-High & Medium & Do if data coverage is acceptable \\
IV / 2SLS full branch & Very High & High (if valid) & Medium-Low & Do only with credible instrument story \\
Dynamic panel & Medium-High & Medium & Low-Medium & Optional, only after FE core is stable \\
Robustness pack (definitions + outliers + splits) & High & Medium-High & High & Strong 14-day candidate \\
Reproducible package + one-pager & Very High & N/A & Very High & Mandatory for portfolio value \\
Forecasting extension & Medium & Low-Medium & Medium & Secondary branch \\
Policy/risk memo & High & N/A & High & Strong add-on after core results \\
\bottomrule
\end{tabular}
\end{center}

\section{Practical Continue / Stop Triggers}
\textbf{Continue} if within first 3 work sessions you can produce:
\begin{itemize}
\item A stable two-way FE result with interpretable sign/magnitude.
\item At least two robustness checks that do not reverse the baseline story completely.
\item One concise methods/results writeup usable in interview discussion.
\end{itemize}

\textbf{Stop or pivot} if any occurs:
\begin{itemize}
\item Instrument idea cannot be defended after quick validity check.
\item Data merges for controls become a multiday bottleneck.
\item Scope grows beyond 14-day cap without clear incremental portfolio gain.
\end{itemize}

\section{Minimal High-Value Path (if continuing)}
This is not a fixed plan, but the highest expected-value branch under your filter:
\begin{enumerate}
\item Two-way FE baseline.
\item Robustness pack (2 to 4 checks).
\item Reproducible package plus one-page narrative.
\item Optional IV add-on only if instrument quality survives fast diagnostics.
\end{enumerate}

\section{Open Idea Backlog (Uncommitted)}
\begin{itemize}
\item Heterogeneous effects by inflation regime.
\item Event-window around major monetary regime changes.
\item Country-clustered standard errors vs alternative clustering structures.
\item Bayesian shrinkage variant for noisy country-level estimates.
\item Alternative money aggregates if data quality permits.
\end{itemize}

\end{document}
